\section{Problem Context}

Locating business rules in large-scale software systems is a well-known challenge in industry. In many organizations, business rules are deeply embedded in legacy code and are not explicitly linked to their descriptions. As a result, identifying the code that implements a particular rule can be difficult, especially when the rule needs to be changed because of maintenance activities or regulatory updates.

In this work, we study a concrete industrial case provided by Berger-Levrault, a major international software company that develops information systems for public administration, healthcare, and education. Like many long-established software companies, Berger-Levrault has large and complex legacy systems in which business rules are closely intertwined with application logic. As these rules evolve, identifying and modifying their implementations can therefore become a challenging task for developers.

When a business rule changes, developers first need to determine which parts of the source code implement that rule. In a large and continuously evolving codebase, identifying the relevant classes and methods is not straightforward. Developers may need to inspect and navigate a large number of files before finding the code related to the rule. This makes maintenance activities more time-consuming and increases the risk of modifying code that is not directly related to the rule, potentially introducing unintended side effects.

Although documentation is available, it is incomplete and may be outdated, which is common in long-lived software systems. It does not provide a complete view of the business rules implemented in the system, nor does it explicitly link the textual description of a rule to its implementation in the source code. Therefore, documentation alone cannot provide reliable support for maintenance tasks, particularly when business regulations or organizational requirements change.

In practice, the source code has consequently become the main source of information for understanding how business rules are implemented. However, the codebase is large, has evolved over many years, and has been maintained and extended by different development teams. Manually navigating such a codebase to locate the implementation of a specific business rule is therefore time-consuming and error-prone.

This industrial context highlights the need for automated approaches that can help developers locate business rules directly in the source code. Such approaches can reduce the effort required to identify relevant implementations and support maintenance activities without relying solely on documentation or external specifications.